\documentclass[11pt]{beamer}

% Idioma y codificación
\usepackage[utf8]{inputenc}
\usepackage[T1]{fontenc}
\usepackage[spanish]{babel}
\usepackage{booktabs}
\usepackage{tabularx}
\usepackage{graphicx}
\usepackage{xcolor}

% Tema y colores personalizados
\usetheme{Madrid}
\usecolortheme{whale}

% Colores personalizados tipo Odín
\definecolor{odinDarkBlue}{RGB}{26, 54, 93}
\definecolor{odinLightBlue}{RGB}{66, 153, 225}
\definecolor{odinBg}{RGB}{247, 250, 252}

\setbeamercolor{palette primary}{bg=odinDarkBlue, fg=white}
\setbeamercolor{palette secondary}{bg=odinLightBlue, fg=white}
\setbeamercolor{palette tertiary}{bg=odinDarkBlue, fg=white}
\setbeamercolor{titlelike}{parent=palette primary}
\setbeamercolor{structure}{fg=odinDarkBlue}

% Quitar barra de navegación inferior molesta
\setbeamertemplate{navigation symbols}{}

% Datos de la presentación
\title[Proyecto Odín]{Proyecto Odín: el asistente total}
\subtitle{Asistente personal local-first, soberano y proactivo}
\author[Enrique de Vicente-Tutor]{Enrique de Vicente-Tutor Castillo}
\institute[FIB - UPC]{
  Facultat d'Informàtica de Barcelona (FIB)\\
  Universitat Politècnica de Catalunya (UPC) - BarcelonaTech
}
\date[Curso 2025-2026]{Defensa de Trabajo de Fin de Grado}

\begin{document}

% Diapositiva 1: Portada
\begin{frame}
  \titlepage
\end{frame}

% Diapositiva 2: Agenda
\begin{frame}{Agenda}
  \tableofcontents
\end{frame}

\section{Introducción y Objetivos}

% Diapositiva 3: Introducción y Motivación
\begin{frame}{Introducción y Motivación}
  \begin{block}{El problema de los asistentes comerciales (Alexa, Siri, Google Home)}
    \begin{itemize}
      \item Dependencia absoluta de la nube y pérdida de privacidad.
      \item Respuestas genéricas y falta de personalización profunda.
      \item Vulnerabilidad ante caídas de internet y obsolescencia de hardware.
    \end{itemize}
  \end{block}

  \begin{block}{La Solución: Odín}
    \begin{itemize}
      \item **Local-first**: inferencia, almacenamiento y procesamiento en local.
      \item **Soberano**: el usuario tiene el control absoluto de sus datos.
      \item **Proactivo**: no solo reacciona, sino que monitoriza, alerta y se autorepara.
    \end{itemize}
  \end{block}
\end{frame}

% Diapositiva 4: Objetivos del Proyecto
\begin{frame}{Objetivos del Proyecto}
  \begin{columns}
    \begin{column}{0.5\textwidth}
      \begin{block}{Objetivo Principal}
        Diseñar, desplegar y validar un ecosistema autoalojado de asistencia personal con IA local y domótica bajo principios de privacidad por diseño.
      \end{block}
    \end{column}
    \begin{column}{0.5\textwidth}
      \begin{block}{Objetivos Secundarios}
        \begin{itemize}
          \item Inferencia LLM local viable.
          \item Memoria documental RAG.
          \item Control domótico (Home Assistant).
          \item Acceso remoto seguro (Tailscale).
          \item Monitorización y autoreparación.
        \end{itemize}
      \end{block}
    \end{column}
  \end{columns}
\end{frame}

\section{Diseño e Infraestructura}

% Diapositiva 5: Arquitectura del Sistema
\begin{frame}{Arquitectura del Sistema (Capas)}
  \begin{table}
    \centering
    \caption{Capas de la arquitectura de Odín}
    \label{tab:capas}
    \begin{tabularx}{\textwidth}{lX}
      \toprule
      Capa & Componentes principales \\
      \midrule
      \textbf{Interacción} & Open WebUI, HA Voice, Telegram, Dashboard \\
      \textbf{Cognición (IA)} & Ollama (\texttt{ministral-3:14b}), Wyoming, Piper, Faster Whisper \\
      \textbf{Datos y Memoria} & Nextcloud, Qdrant (vectorial), Immich, Mealie \\
      \textbf{Físico y Domótico} & Servidor principal, Raspberry Pi 5, Home Assistant \\
      \textbf{Acceso y Red} & Tailscale VPN, proxies locales y Caddy \\
      \bottomrule
    \end{tabularx}
  \end{table}
\end{frame}

% Diapositiva 6: Despliegue de Dos Nodos
\begin{frame}{Distribución de Nodos: Servidor y Raspberry Pi 5}
  \begin{block}{Nodo 1: Servidor Principal (CPU/GPU AMD)}
    \begin{itemize}
      \item Concentra la carga de cómputo pesado: LLM (Ollama), base vectorial (Qdrant), fotos (Immich), nube (Nextcloud) y orquestación (n8n).
      \item Despliegue reproducible mediante contenedores Docker Compose.
    \end{itemize}
  \end{block}
  \begin{block}{Nodo 2: Raspberry Pi 5 (Home Assistant)}
    \begin{itemize}
      \item Ejecuta Home Assistant OS de forma nativa e independiente.
      \item Garantiza la alta disponibilidad de la domótica física aunque el servidor se apague o esté en mantenimiento.
    \end{itemize}
  \end{block}
\end{frame}

\section{Desarrollo e Integración}

% Diapositiva 7: Interacción Conversacional y Tools
\begin{frame}{Interacción Conversacional y Tools}
  \begin{block}{Open WebUI como Interfaz Principal}
    \begin{itemize}
      \item Conectado al modelo local \texttt{ministral-3:14b} con un contexto optimizado de \textbf{16,384 tokens} para evitar truncados en lecturas largas.
    \end{itemize}
  \end{block}
  \begin{block}{El Catálogo de Tools (Acciones Atómicas)}
    El modelo decide cuándo usar herramientas para interactuar con el mundo real:
    \begin{itemize}
      \item \textbf{Home Assistant}: control de luces, clima, calendario y aspirador Dreame.
      \item \textbf{Memoria Vectorial}: búsquedas semánticas y guardado de notas.
      \item \textbf{Immich / Frigate}: recuperación de fotos y últimas alertas de cámara.
      \item \textbf{Mealie}: creación e importación directa de recetas.
    \end{itemize}
  \end{block}
\end{frame}

% Diapositiva 8: Proxies Locales Especializados
\begin{frame}{Proxies Locales para Seguridad y CORS}
  \begin{itemize}
    \item Para evitar exponer API keys y solucionar problemas de CORS al renderizar recursos privados en el navegador, se crearon proxies específicos en contenedores:
  \end{itemize}
  \begin{table}
    \centering
    \begin{tabular}{lcl}
      \toprule
      Contenedor Proxy & Puerto & Finalidad \\
      \midrule
      \texttt{immich-photo-proxy} & 2290 & Renderiza imágenes en Open WebUI. \\
      \texttt{nextcloud-note-proxy} & 2291 & Envía notas de forma robusta a n8n. \\
      \texttt{mealie-tool-proxy} & 2292 & Permite importar recetas sin exponer tokens. \\
      \texttt{frigate-image-proxy} & 2293 & Sirve capturas de cámaras con CORS. \\
      \bottomrule
    \end{tabular}
  \end{table}
\end{frame}

\section{Resultados y Pruebas}

% Diapositiva 9: Pruebas Funcionales
\begin{frame}{Pruebas Funcionales y Validación}
  \begin{itemize}
    \item \textbf{Ingesta de Memoria}: Validación horaria con cron. 35 archivos procesados a embeddings en Qdrant (colección \texttt{memoria\_ia} en estado \texttt{green} con 3,554 puntos).
    \item \textbf{Cámaras}: Recuperación exitosa de eventos de personas en Frigate y visualización en chat a través del proxy.
    \item \textbf{Dashboard Domótico}: Integración completa de un panel unificado en Home Assistant (6 vistas, 31 secciones y 69 entidades) sin dependencias externas.
  \end{itemize}
\end{frame}

% Diapositiva 10: Medición de Consumo Real
\begin{frame}{Evidencia Experimental: Medición de Consumo Real}
  \begin{block}{Resultados de los dos Enchufes Inteligentes Integrados}
    \begin{itemize}
      \item \textbf{Servidor Principal (Torre)}:
        \begin{itemize}
          \item Potencia media registrada: \textbf{109.93 W} (reposo y cargas LLM).
          \item Consumo promedio diario: \textbf{0.64 kWh}.
          \item Gasto mensual proyectado: \textbf{2.88 EUR/mes} (tarifa de 0.15 EUR/kWh).
        \end{itemize}
      \item \textbf{Equipamiento Auxiliar (Raspberry Pi 5 + Red)}:
        \begin{itemize}
          \item Potencia registrada: \textbf{3.81 W}.
          \item Gasto mensual proyectado: \textbf{0.18 EUR/mes}.
        \end{itemize}
      \item \textbf{Total Ecosistema 24/7}: \textbf{113.74 W} de media y un coste de \textbf{3.06 EUR/mes}.
    \end{itemize}
  \end{block}
\end{frame}

\section{Conclusiones y Trabajo Futuro}

% Diapositiva 11: Conclusiones
\begin{frame}{Conclusiones}
  \begin{block}{Logros}
    \begin{itemize}
      \item Viabilidad demostrada de un asistente local-first moderno y privado.
      \item Arquitectura modular robusta y mantenible (Docker + Home Assistant).
      \item Integración de telemetría y alertas proactivas (Telegram, Netdata).
    \end{itemize}
  \end{block}
  \begin{block}{Limitaciones Aceptadas}
    \begin{itemize}
      \item La voz interactiva (Piper/ASR) funciona pero carece de la velocidad y naturalidad de la nube comercial.
      \item El backup físico con NVMe externo está en fase de integración final.
    \end{itemize}
  \end{block}
\end{frame}

% Diapositiva 12: Trabajo Futuro
\begin{frame}{Trabajo Futuro}
  \begin{itemize}
    \item \textbf{Consolidación del Backup Físico}: Automatizar el montaje por UUID del NVMe, cifrado de datos y pruebas de restauración.
    \item \textbf{Estudio Longitudinal del Consumo}: Evaluar variaciones estacionales y picos durante cargas intensivas continuas de inferencia.
    \item \textbf{Asistente de Escritorio}: Extender Odín para que pueda actuar sobre el PC del usuario (gestión de archivos locales, ejecución asistida de comandos con confirmación).
  \end{itemize}
  \begin{center}
    \Huge \bfseries ¡Muchas gracias!\\
    \large ¿Preguntas?
  \end{center}
\end{frame}

\end{document}
